\chapter{Invariants and Exploitability}

\section{Purpose}

The previous chapter ended with a demand:

\begin{quote}
A threat model must produce invariants, tests, mitigations, monitoring,
and residual-risk decisions.
\end{quote}

This chapter teaches the first and most important output: invariants.

An invariant is not a slogan like ``funds should be safe.'' It is a
property that should remain true across the states, transitions, or
histories the system can actually reach. In blockchain security, the
deepest bugs are often not parser bugs, syntax errors, or obviously
invalid transactions. They are valid sequences of transactions that move
the system into a state the protocol designers did not mean to allow.

That is the central lesson of Part I:

\begin{lstlisting}
An exploit is often a counterexample to an intended invariant.
\end{lstlisting}

If the invariant is vague, the review is vague. If the invariant is
wrong, the tests will prove the wrong thing. If the invariant ignores
capital, timing, ordering, liquidity, privilege, and profit, it may
describe a theoretical weakness but not a realistic exploit. This
chapter connects those pieces.

By the end of this chapter, a reader should be able to look at a
protocol description and extract concrete security properties such as:

\begin{itemize}
\tightlist
\item
  Total issued shares must correspond to a claim on real assets.
\item
  No account's collateral may be seized unless a liquidation condition
  is true.
\item
  A bridge must not mint more wrapped assets than are locked or
  otherwise secured on the source domain.
\item
  A withdrawal queue must not allow later users to bypass earlier users
  if priority is part of the design.
\item
  A vault donation must not let an early depositor steal a later
  depositor's assets through rounding.
\item
  A governance payload must not execute before the required timelock has
  elapsed.
\end{itemize}

Then the reader must go further:

\begin{itemize}
\tightlist
\item
  Can an attacker reach a violating state?
\item
  What sequence gets there?
\item
  How much capital is needed?
\item
  Does the attacker need ordering control?
\item
  Is the attack profitable after fees, slippage, gas, liquidity limits,
  and competition?
\item
  Does the attack break safety, liveness, solvency, fairness, authority,
  or user expectation?
\end{itemize}

That second set of questions is exploitability.

\section{Learning Objectives}

After studying this chapter, you should be able to:

\begin{itemize}
\tightlist
\item
  Define an invariant as a property over reachable states, transitions,
  or histories.
\item
  Distinguish invariants from aspirations, tests, assertions,
  assumptions, and implementation details.
\item
  Write conservation, authorization, solvency, temporal, and economic
  invariants for blockchain systems.
\item
  Explain why many attacks are valid state transitions that violate
  protocol intent.
\item
  Translate a broken invariant into an attack hypothesis.
\item
  Evaluate exploitability using capital, timing, ordering, liquidity,
  reliability, privilege, and profit.
\item
  Recognize weak invariants that are too broad, too narrow, untestable,
  circular, or assumption-heavy.
\item
  Use invariant thinking as the bridge from threat modeling to testing,
  fuzzing, formal verification, and manual review.
\end{itemize}

\section{The Core Model}

For Part I, use this mental model:

\begin{lstlisting}
System:
  state space S
  initial states I
  transition relation T
  assumptions A

Invariant:
  property P that should hold over all reachable states, transitions, or histories

Security failure:
  a reachable state or history that violates P under assumptions A

Exploit:
  a feasible strategy that drives the system to such a violation and extracts value,
  control, denial, or other advantage
\end{lstlisting}

For a simple state invariant:

\begin{lstlisting}
P holds initially.
Every allowed transition preserves P.
Therefore P holds in every reachable state.
\end{lstlisting}

This is the classic inductive invariant method.

Formal notation can be useful, but keep it tied to review work:

\begin{lstlisting}
For all reachable states S:
  P(S) == true

For every allowed transition T:
  P(S) == true implies P(apply(S, T)) == true

Counterexample:
  find S and T such that P(S) == true but P(apply(S, T)) == false
\end{lstlisting}

The notation is not decoration. It tells you what kind of evidence would falsify the claim.
 Lamport describes state
machines in terms of initial states and next-state relations, then
treats invariance as a correctness method over reachable
states.\footnote{Leslie Lamport, ``Computation and State Machines'',
  https://lamport.azurewebsites.net/pubs/state-machine.pdf.} Certora's
documentation makes the same proof structure practical for smart
contracts: establish the invariant after construction, then check that
every relevant method preserves it if it held before the
method.\footnote{Certora Prover Documentation, ``Invariants'', invariant
  establishment and preservation checks,
  https://docs.certora.com/en/latest/docs/cvl/invariants.html.}

In security review, however, the practical question is usually inverted:

\begin{lstlisting}
Can I find a valid transition sequence from an initial state to a state where P is false?
\end{lstlisting}

That sequence is a counterexample. In blockchain, a counterexample often
looks like an exploit trace.

\section{What an Invariant Is and Why It Finds Deep Bugs}

\subsection{Invariants Are Properties, Not Wishes}

A weak statement:

\begin{lstlisting}
Users should not lose funds.
\end{lstlisting}

A better invariant:

\begin{lstlisting}
For every user u, after a successful withdrawal of s shares,
the assets transferred to u must be no less than previewRedeem(s)
computed immediately before the withdrawal, except for documented fees.
\end{lstlisting}

Another weak statement:

\begin{lstlisting}
Only governance can upgrade the protocol.
\end{lstlisting}

A better invariant:

\begin{lstlisting}
No implementation address used by the proxy can change unless the call path
originates from the governance executor after the proposal has passed,
the timelock delay has elapsed, and the payload target and calldata match
the queued proposal.
\end{lstlisting}

Another weak statement:

\begin{lstlisting}
The bridge should be backed.
\end{lstlisting}

A better invariant:

\begin{lstlisting}
For asset a, total wrapped supply on destination domains must be less than or
equal to source-domain locked balance plus explicitly recognized canonical debt,
minus completed burns and withdrawals, under the bridge's finality assumptions.
\end{lstlisting}

The difference is not style. It is reviewability. A wish cannot be
disproved. A useful invariant can be attacked.

\subsection{Reachability Is the Point}

An invariant is not merely a true statement about the system at rest. It
must hold over reachable states.

This matters because many bad reviews inspect only clean, intended
states:

\begin{itemize}
\tightlist
\item
  The vault has normal deposits.
\item
  The oracle price is fresh.
\item
  The pool has healthy liquidity.
\item
  The bridge validator set is honest.
\item
  The admin key acts according to policy.
\item
  The contract is called through the official frontend.
\end{itemize}

Attackers search the rest of the reachable state space:

\begin{itemize}
\tightlist
\item
  First deposit is one wei.
\item
  A token charges transfer fees.
\item
  A donation changes \passthrough{\lstinline!balanceOf!} without
  updating accounting.
\item
  A borrow occurs immediately after an oracle update.
\item
  A reentrant callback observes state before restoration.
\item
  A liquidation races a price change.
\item
  A bridge message is replayed across domains.
\item
  A governance proposal executes through an unexpected call path.
\end{itemize}

Invariant thinking forces a reviewer to ask:

\begin{lstlisting}
What can be true before this transaction?
What can this transaction change?
What must still be true after it?
\end{lstlisting}

\subsection{Inductive Invariants vs Intended Invariants}

An invariant can be true of the implementation and still be useless.

Example:

\begin{lstlisting}
totalSupply is always equal to totalSupply.
\end{lstlisting}

This is true and worthless.

Another example:

\begin{lstlisting}
The contract's stored `totalDeposits` variable never decreases except in withdraw.
\end{lstlisting}

This may be true but incomplete if token balances can change through
direct transfer, rebasing, fee-on-transfer behavior, slashing, or
strategy loss.

Security review cares about intended invariants:

\begin{lstlisting}
The accounting variables that define user claims must correspond to actual
recoverable assets under the token and strategy semantics the protocol supports.
\end{lstlisting}

The gap between the implemented invariant and the intended invariant is
where deep bugs live.

\subsection{Invariants Can Be About States, Transitions, or Histories}

Most beginners think invariants are only static state predicates:

\begin{lstlisting}
assetBalance >= totalUserClaims
\end{lstlisting}

Those are important, but not enough.

Blockchain systems need several kinds of properties.

{\def\LTcaptype{none} % do not increment counter
\begin{longtable}[]{@{}
  >{\raggedright\arraybackslash}p{(\linewidth - 4\tabcolsep) * \real{0.3333}}
  >{\raggedright\arraybackslash}p{(\linewidth - 4\tabcolsep) * \real{0.3333}}
  >{\raggedright\arraybackslash}p{(\linewidth - 4\tabcolsep) * \real{0.3333}}@{}}
\toprule\noalign{}
\begin{minipage}[b]{\linewidth}\raggedright
Property type
\end{minipage} & \begin{minipage}[b]{\linewidth}\raggedright
Question
\end{minipage} & \begin{minipage}[b]{\linewidth}\raggedright
Example
\end{minipage} \\
\midrule\noalign{}
\endhead
\bottomrule\noalign{}
\endlastfoot
State invariant & What must be true now? &
\passthrough{\lstinline!cash + borrows - reserves >= liabilities!} \\
Transition invariant & What relationship must hold before and after one
action? & A withdrawal burns shares before reducing total assets owed \\
Relational invariant & What variables must agree? & Sum of per-user
balances equals total supply \\
Temporal property & What must eventually or never happen over time? & A
queued withdrawal eventually becomes claimable if the system is
unpaused \\
Ordering property & What must happen before what? & A governance payload
cannot execute before its timelock \\
Economic invariant & What value relationship must hold? & A solvent
account cannot be liquidated at the current oracle price \\
Observational invariant & What must external observers be able to
verify? & A bridge proof must bind source chain, message nonce, sender,
recipient, and amount \\
\end{longtable}
}

The safety/liveness distinction is useful here. Safety says ``bad things
do not happen''; liveness says ``good things eventually happen.'' Alpern
and Schneider's classical work gives the formal background for that
split.\footnote{Bowen Alpern and Fred B. Schneider, ``Defining
  Liveness'',
  https://www.cs.cornell.edu/fbs/publications/DefLiveness.pdf.} In
blockchain security, most theft bugs are safety failures, while stuck
withdrawals, unprocessable queues, permanent censorship, and unusable
exits are often liveness failures.

Do not force every property into a state invariant. Use the right shape.

\subsection{A Broken Invariant Is an Attack Hypothesis}

When you write:

\begin{lstlisting}
For every account, debtValue <= collateralValue * collateralFactor
after any borrow or withdrawal.
\end{lstlisting}

you are implicitly asking:

\begin{lstlisting}
Can an attacker make debtValue exceed the allowed collateralized value?
\end{lstlisting}

That becomes an attack hypothesis:

\begin{lstlisting}
If the attacker can manipulate the oracle price used in collateralValue,
then the borrow check may pass at an inflated collateral value.
After price normalization, the account may be insolvent.
\end{lstlisting}

When you write:

\begin{lstlisting}
Only the owner of a position can withdraw its collateral,
unless the position is liquidatable.
\end{lstlisting}

you are asking:

\begin{lstlisting}
Can the attacker make the system treat them as owner,
or make a healthy position appear liquidatable,
or call a path that bypasses this check?
\end{lstlisting}

This is why invariants find deep bugs. They do not ask whether a known
vulnerability label appears. They ask whether the system can violate its
own security claims.

\subsection{Tools Check Properties; People Must Choose Them}

Invariant testing and formal tools are valuable, but they do not decide
what the protocol should mean.

Foundry describes invariant testing as checking properties that should
always hold regardless of the sequence of actions, with Forge generating
random call sequences and checking invariants after calls.\footnote{Foundry,
  ``Invariant Testing'',
  https://www.getfoundry.sh/guides/invariant-testing.} Echidna similarly
generates arbitrary transactions and reports sequences that make a
property fail.\footnote{Trail of Bits, ``Testing a Property with
  Echidna'',
  https://secure-contracts.com/program-analysis/echidna/introduction/how-to-test-a-property.html.}
Solidity's SMTChecker documentation is blunt about the specification
problem: formal verification helps compare what you specified with what
you implemented, but you still must check whether the specification is
what you wanted.\footnote{Solidity documentation, ``SMTChecker and
  Formal Verification'',
  https://docs.soliditylang.org/en/latest/smtchecker.html.}

That is the reviewer's burden.

Bad invariant:

\begin{lstlisting}
function invariant_totalAssetsNonzero() external view {
    assert(vault.totalAssets() >= 0);
}
\end{lstlisting}

This proves nothing useful.

Better invariant:

\begin{lstlisting}
function invariant_solvency() external view {
    assert(asset.balanceOf(address(vault)) + strategyRecoverableAssets()
        >= vault.totalAssetsOwedToShareholders());
}
\end{lstlisting}

Still incomplete unless:

\begin{itemize}
\tightlist
\item
  \passthrough{\lstinline!strategyRecoverableAssets()!} is defensible.
\item
  The asset token behavior is modeled correctly.
\item
  All share holders and debt-like claims are included.
\item
  Fees, pending withdrawals, locked profit, slashing, and bad debt are
  accounted for.
\item
  The invariant is checked across reachable sequences, not only a single
  happy path.
\end{itemize}

Tools amplify good properties. They also amplify bad ones.

\subsection{The Invariant Quality Test}

Before using an invariant in review, testing, fuzzing, or formal
verification, test the invariant itself.

A useful invariant is:

\begin{itemize}
\tightlist
\item
  \textbf{Specific:} It names the state, asset, actor, authority, or
  transition it constrains.
\item
  \textbf{Falsifiable:} A counterexample can be described.
\item
  \textbf{Reachability-aware:} It applies to states the system can
  actually reach, including adversarial sequences.
\item
  \textbf{Assumption-aware:} It states which oracle, finality, token,
  governance, or operator assumptions it needs.
\item
  \textbf{Quantified:} It says whether it applies to all users, all
  assets, all markets, all messages, or a subset.
\item
  \textbf{Measurable:} A reviewer or tool can evaluate it from code,
  state, traces, or explicit models.
\item
  \textbf{Non-circular:} It does not merely restate the implementation
  variable as truth.
\item
  \textbf{Not too weak:} Passing it would actually reduce risk.
\item
  \textbf{Not too strong:} It does not forbid intended behavior.
\item
  \textbf{Linked to impact:} Violating it would matter.
\end{itemize}

If an invariant fails this test, refine it before trusting the review
built on it.

\section{Conservation, Supply, and Asset-Balance Invariants}

\subsection{Conservation Is the First Accounting Question}

For systems that mint, burn, lock, unlock, wrap, deposit, withdraw,
swap, stake, or bridge assets, the first security question is:

\begin{lstlisting}
Where did the value come from, and where can it go?
\end{lstlisting}

Conservation invariants track value across representations.

Examples:

\begin{lstlisting}
Token:
  totalSupply == sum(balanceOf(account) for all accounts)

Vault:
  totalAssetsManaged >= assetsRedeemableByAllShares

Bridge:
  wrappedSupply(destination) <= lockedAssets(source) + recognizedCredit - pendingExits

AMM:
  reserves after swap must satisfy the pool's pricing invariant after fees

Reward distributor:
  totalClaimed + totalUnclaimed + undistributedBalance <= totalFundedRewards
\end{lstlisting}

These are not all the same kind of invariant. Some are exact. Some are
inequalities. Some depend on external facts. Some need a model because
summing all accounts on-chain may not be feasible. But the underlying
question is the same:

\begin{lstlisting}
Can claims exceed backing?
\end{lstlisting}

\subsection{Exact Conservation vs Bounded Drift}

Do not demand exact equality when the system intentionally permits
drift.

Exact:

\begin{lstlisting}
ERC-20 totalSupply equals the sum of all balances.
\end{lstlisting}

Bounded:

\begin{lstlisting}
Total rounding loss allocated to users must be less than one unit per operation
and must not be accumulable by an attacker into a net profit above fees.
\end{lstlisting}

Intentional inequality:

\begin{lstlisting}
Vault assets may exceed shareholder claims because donations or strategy yield
can increase assets without minting shares.
\end{lstlisting}

Dangerous inequality:

\begin{lstlisting}
Vault assets are less than shareholder claims because strategy losses were not
recognized or withdrawal accounting double-counted assets.
\end{lstlisting}

Many accounting bugs come from confusing these cases.

\subsection{Total Supply Is Not Enough}

Beginners often stop at:

\begin{lstlisting}
totalSupply == sum(balances)
\end{lstlisting}

That is necessary for many tokens, but it does not establish asset
safety.

A wrapped token can have perfect internal supply accounting and still be
unbacked. A lending market can track cToken balances correctly and still
become insolvent. A vault can correctly sum shares and still issue
shares at the wrong exchange rate. A bridge can correctly burn wrapped
tokens and still release source assets to the wrong recipient.

Use layers:

\begin{lstlisting}
Internal accounting:
  totalSupply == sum(userShares)

Claim accounting:
  userAssetsClaim = sharesToAssets(userShares)

Backing:
  recoverableAssets >= sum(userAssetsClaim)

Authority:
  only permitted transitions can change supply, claims, or backing

External dependency:
  recoverableAssets calculation is valid under token, oracle, strategy, and bridge assumptions
\end{lstlisting}

The strongest bugs often appear between layers.

\subsection{Direct Balance Changes Are a Security Surface}

Many protocols use token balances as accounting inputs:

\begin{lstlisting}
assets = asset.balanceOf(address(this))
\end{lstlisting}

That may be correct. It may also be dangerous.

Token balances can change without the protocol's deposit function:

\begin{itemize}
\tightlist
\item
  Direct transfers.
\item
  Rebases.
\item
  Fee-on-transfer behavior.
\item
  Token hooks.
\item
  Token recovery functions.
\item
  Strategy returns.
\item
  Airdrops.
\item
  Slashing.
\item
  Cross-chain minting or burning.
\end{itemize}

If a protocol's accounting assumes all asset changes pass through its
own functions, direct balance changes break the model.

Example invariant:

\begin{lstlisting}
If share price is computed from asset.balanceOf(vault),
then direct asset transfers to the vault must not allow an attacker to extract
more value from later users than the attacker contributed.
\end{lstlisting}

That is a real security claim, not a generic ``handle donations'' note.

\subsection{Bridges Need Cross-Domain Conservation}

A lock-mint bridge has a simple intended accounting story:

\begin{lstlisting}
source chain: lock original token
destination chain: mint wrapped token
\end{lstlisting}

The invariant is not local to either chain:

\begin{lstlisting}
wrappedSupply(destination) <= lockedBalance(source)
\end{lstlisting}

But real bridges complicate this:

\begin{itemize}
\tightlist
\item
  Multiple destination chains.
\item
  Multiple token representations.
\item
  Pending messages.
\item
  Failed withdrawals.
\item
  Finality windows.
\item
  Validator signatures.
\item
  Relayer retries.
\item
  Fee deductions.
\item
  Emergency pauses.
\item
  Canonical debt during asynchronous settlement.
\end{itemize}

So the real invariant becomes:

\begin{lstlisting}
For each asset a:
  sum(wrappedSupply[a] across destination domains)
  + pendingUnlocks[a]
  - pendingBurns[a]
  <= sourceLocked[a] + explicitlyAuthorizedCredit[a]

under the bridge's finality and message-validity assumptions.
\end{lstlisting}

The assumption clause matters. A bridge invariant may be true if
validators never sign fraudulent messages and false if a quorum can be
compromised. That does not make the invariant useless. It tells you
where the security boundary is.

\subsection{Conservation Invariants Should Name Units}

Always name units.

Bad:

\begin{lstlisting}
balances match.
\end{lstlisting}

Better:

\begin{lstlisting}
The sum of user principal balances denominated in underlying asset units
must be less than or equal to recoverable underlying assets.
\end{lstlisting}

Better still:

\begin{lstlisting}
For USDC market m, after applying USDC's 6-decimal precision and the protocol's
index scale, the sum of user supply claims in USDC atomic units must be less
than or equal to cash + borrows - reserves - recognizedBadDebt, allowing at most
1 atomic unit of rounding error per account update.
\end{lstlisting}

Units are not clerical. Decimal mismatch, index scale mismatch, and
rounding direction are exploit surfaces.

\subsection{Common Conservation Failure Modes}

{\def\LTcaptype{none} % do not increment counter
\begin{longtable}[]{@{}
  >{\raggedright\arraybackslash}p{(\linewidth - 2\tabcolsep) * \real{0.5000}}
  >{\raggedright\arraybackslash}p{(\linewidth - 2\tabcolsep) * \real{0.5000}}@{}}
\toprule\noalign{}
\begin{minipage}[b]{\linewidth}\raggedright
Failure mode
\end{minipage} & \begin{minipage}[b]{\linewidth}\raggedright
Broken invariant
\end{minipage} \\
\midrule\noalign{}
\endhead
\bottomrule\noalign{}
\endlastfoot
Mint without backing & Supply exceeds assets, locks, or authorized
credit \\
Burn without releasing claim & User loses claim without receiving assets
or debt reduction \\
Double claim & Same deposit creates multiple redeemable claims \\
Untracked liability & Pending withdrawals or rewards omitted from
solvency calculation \\
Donation manipulation & Direct balance change shifts exchange rate in
attacker's favor \\
Fee-on-transfer mismatch & Protocol credits nominal amount but receives
less \\
Rebase mismatch & Internal accounting assumes balances are stable \\
Cross-domain replay & Same source event mints or unlocks more than
once \\
Decimal mismatch & Value comparisons mix incompatible units \\
Rounding accumulation & Tiny errors become extractable over repeated
calls \\
\end{longtable}
}

For every asset-moving feature, write at least one conservation
invariant before reading the implementation.

\section{Authorization, Ownership, and Capability Invariants}

\subsection{Authorization Is a State Property}

Authorization is often taught as a function-level check:

\begin{lstlisting}
require(msg.sender == owner);
\end{lstlisting}

That is too small. Authorization is a property of the state transition:

\begin{lstlisting}
Who is allowed to cause this state change, under what conditions, through which path?
\end{lstlisting}

A better invariant:

\begin{lstlisting}
No transition may reduce user u's asset claim unless one of the following is true:
  u authorized the transition,
  u's position satisfies the liquidation predicate,
  a documented fee was charged according to protocol rules,
  or an explicitly scoped emergency process executed.
\end{lstlisting}

This captures more than \passthrough{\lstinline!msg.sender!}. It
includes signatures, approvals, liquidations, fee logic, governance, and
emergency controls.

\subsection{Ownership Must Be Bound to the Right Object}

Ownership bugs are often object-binding bugs.

Examples:

\begin{itemize}
\tightlist
\item
  Checking that the caller owns some NFT, but not the NFT controlling
  this position.
\item
  Checking that an account is a signer, but not that it signed this
  message.
\item
  Checking that a PDA has the expected derivation, but not that it
  belongs to the expected program and seed domain.
\item
  Checking that a proof is valid, but not that it proves inclusion for
  this chain, token, amount, recipient, and nonce.
\item
  Checking that a governance proposal passed, but not that the executed
  calldata matches the proposal.
\end{itemize}

Authorization invariant:

\begin{lstlisting}
For every protected object o and protected action a,
the authority used to authorize a must be bound to o, a, the target contract,
the chain or domain, and any nonce or validity window required to prevent replay.
\end{lstlisting}

This repeats a lesson from Chapter 5: signatures authorize messages, not
intentions. The same is true for ownership checks. They authorize
exactly what they bind.

\subsection{Capabilities Are Assets}

Some systems do not use role strings or owner addresses as the main
authority model. They use capabilities:

\begin{itemize}
\tightlist
\item
  A Move resource that grants mint authority.
\item
  A Solana PDA that controls token accounts.
\item
  A smart account session key with scoped permissions.
\item
  A governance executor contract that can call privileged targets.
\item
  A multisig module that can move treasury funds.
\item
  A validator key that can sign bridge messages.
\end{itemize}

A capability is an asset because controlling it changes what state
transitions are possible.

Capability invariant:

\begin{lstlisting}
A capability that can mint, burn, upgrade, pause, seize, withdraw, or sign
cross-domain messages must only be created, transferred, delegated, or revoked
through explicitly authorized transitions.
\end{lstlisting}

Then specialize:

\begin{lstlisting}
Mint capability:
  Only governance can grant it.
  Granting requires timelock.
  Revoked minters cannot mint.
  The zero address cannot receive it.
  A minter cannot exceed per-epoch limits.

Session key:
  It cannot spend outside its token allowlist.
  It cannot exceed its value cap.
  It expires at or before its configured deadline.
  It cannot upgrade its own permissions.

Bridge signer:
  It cannot sign messages for unsupported domains.
  Its key rotation must not reduce threshold safety below the configured quorum.
\end{lstlisting}

Do not treat capabilities as ``configuration.'' They are part of the
asset map.

\subsection{Negative Authorization Matters}

Authorization is not only ``who can do X.'' It is also ``who cannot do
X.''

Examples:

\begin{lstlisting}
Liquidator:
  Can seize collateral only if liquidation predicate is true.
  Cannot choose arbitrary collateral unrelated to the unhealthy debt.
  Cannot seize more than close factor and incentive allow.

Guardian:
  Can pause deposits and borrows.
  Cannot withdraw user funds.
  Cannot upgrade implementation.
  Cannot unpause after a governance-defined maximum incident window, unless governance confirms.

Strategist:
  Can move assets into approved strategies.
  Cannot move assets into unapproved tokens.
  Cannot set withdrawal fees.
  Cannot change oracle source.
\end{lstlisting}

A role without negative boundaries is not a role. It is a vague trust
assumption.

\subsection{Cross-Contract Authorization Is Where Bugs Hide}

Many protocols split authority across contracts:

\begin{itemize}
\tightlist
\item
  Proxy admin.
\item
  Implementation contract.
\item
  Timelock.
\item
  Governance executor.
\item
  Access-control registry.
\item
  Vault.
\item
  Strategy.
\item
  Oracle adapter.
\item
  Token.
\item
  Router.
\end{itemize}

Local checks can pass while the system-level authorization invariant
fails.

Example:

\begin{lstlisting}
The strategy checks that caller == vault.
The vault lets anyone call `executeStrategyCall(bytes calldata data)`.
The system-level effect is that anyone can make the strategy perform privileged actions.
\end{lstlisting}

The local invariant ``strategy only trusts vault'' is true. The system
invariant ``only authorized strategists can move strategy funds'' is
false.

System authorization invariant:

\begin{lstlisting}
Every call path that reaches a privileged state change must pass through
the intended authorization predicate, and no public wrapper may make that
predicate trivially true.
\end{lstlisting}

\subsection{Authorization Invariants Should Include Time}

Time changes authority.

Examples:

\begin{itemize}
\tightlist
\item
  A signature expires.
\item
  A session key expires.
\item
  A timelock matures.
\item
  A guardian pause expires.
\item
  A withdrawal delay completes.
\item
  A dispute window closes.
\item
  A validator key rotation becomes active.
\end{itemize}

Temporal authorization invariant:

\begin{lstlisting}
Authority granted for interval [start, end] must not authorize protected
transitions before start or after end.
\end{lstlisting}

Governance example:

\begin{lstlisting}
A queued proposal can execute only if:
  proposal state is succeeded,
  execution time >= queuedTime + timelockDelay,
  execution time <= expirationTime if expiration exists,
  payload hash matches the queued payload,
  and the proposal has not already been executed or canceled.
\end{lstlisting}

If you omit time from authority, replay and premature execution bugs
become easy to miss.

\section{Solvency, Collateralization, Reserve, and Share-Price Invariants}

\subsection{Solvency Is About Claims vs Recoverable Assets}

A protocol is solvent with respect to a claim class when it can satisfy
those claims under the rules and assumptions promised.

Simple:

\begin{lstlisting}
recoverableAssets >= userClaims
\end{lstlisting}

But serious protocols need precision:

\begin{lstlisting}
recoverable underlying assets
  >= immediately withdrawable user claims
  + pending withdrawal claims
  + accrued but unpaid fees owed to users
  + recognized protocol obligations
\end{lstlisting}

Questions:

\begin{itemize}
\tightlist
\item
  Which assets count as recoverable?
\item
  At what price?
\item
  Under what withdrawal delay?
\item
  With what slippage?
\item
  Are strategy positions liquid?
\item
  Are pending losses recognized?
\item
  Are bad debts included?
\item
  Are fees liabilities or protocol equity?
\item
  Can an emergency pause delay payment without violating the stated
  property?
\end{itemize}

Solvency is not a single number. It is a claim about asset
recoverability under stated conditions.

\subsection{Lending Markets Need Collateralization Invariants}

Lending protocols convert assets, prices, collateral factors, interest
indexes, and liquidation rules into a safety condition.

Aave describes health factor as collateral value multiplied by weighted
liquidation threshold divided by borrow value, with health factor below
1 signaling liquidation risk.\footnote{Aave, ``Health Factor \&
  Liquidations'', https://aave.com/help/borrowing/liquidations.}
Compound's Comptroller is its risk management layer, using balances,
prices, and collateral factors to determine liquidity and
shortfall.\footnote{Compound v2, ``Comptroller'',
  https://docs.compound.finance/v2/comptroller/.}

The security invariant behind these systems is:

\begin{lstlisting}
A borrower must not be able to increase debt or withdraw collateral if the
resulting account would violate the protocol's collateralization predicate.
\end{lstlisting}

More concrete:

\begin{lstlisting}
After any borrow, collateral withdrawal, collateral transfer, oracle update
application, liquidation, or interest accrual:

  account is either:
    healthy according to the current risk model,
    liquidatable according to the current risk model,
    or in an explicitly recognized bad-debt/recovery state.
\end{lstlisting}

Do not write only:

\begin{lstlisting}
healthFactor(user) >= 1
\end{lstlisting}

That may be false for accounts that are intentionally liquidatable.
Better:

\begin{lstlisting}
Healthy-only operations may not leave healthFactor(user) < 1.
Liquidation-only operations may execute only when healthFactor(user) < 1
or according to the protocol's documented liquidation eligibility predicate.
\end{lstlisting}

The invariant must distinguish operation classes.

\subsection{Lending Invariants Depend on Oracles}

Collateralization is only as good as the price model.

Bad invariant:

\begin{lstlisting}
debtValue <= collateralValue * collateralFactor
\end{lstlisting}

Better:

\begin{lstlisting}
debtValue(priceSnapshot) <= collateralValue(priceSnapshot) * collateralFactor
where priceSnapshot satisfies the protocol's freshness, source, deviation,
liquidity, and sequencer-availability assumptions.
\end{lstlisting}

Attack hypothesis:

\begin{lstlisting}
If the attacker can manipulate priceSnapshot for one transaction,
then the borrow check may accept an undercollateralized account.
\end{lstlisting}

This is why oracle assumptions must be written into the invariant.
Otherwise, a test may prove the accounting math while ignoring the
exploit.

\subsection{Reserve Invariants Separate User Claims from Protocol
Equity}

Many systems have reserves:

\begin{itemize}
\tightlist
\item
  Lending reserves.
\item
  Insurance funds.
\item
  AMM protocol fees.
\item
  Vault performance fees.
\item
  Liquidation penalties.
\item
  Treasury-owned liquidity.
\end{itemize}

Reserve invariant:

\begin{lstlisting}
Protocol reserves must not be counted simultaneously as user-redeemable backing
and protocol-withdrawable equity.
\end{lstlisting}

Double-counting reserves is a solvency bug.

Example:

\begin{lstlisting}
cash + borrows - reserves >= supplierClaims
\end{lstlisting}

If reserves are omitted:

\begin{lstlisting}
cash + borrows >= supplierClaims
\end{lstlisting}

the protocol may appear solvent while the same assets are promised to
suppliers and the treasury.

\subsection{Share-Price Invariants Are Subtle}

Vaults, staking derivatives, LP tokens, lending receipts, and restaking
positions often use shares.

General model:

\begin{lstlisting}
shares represent a proportional claim on assets.
\end{lstlisting}

ERC-4626 standardizes tokenized vault interfaces such as
\passthrough{\lstinline!totalAssets!},
\passthrough{\lstinline!convertToShares!}, and
\passthrough{\lstinline!convertToAssets!}.\footnote{Ethereum EIP-4626,
  ``Tokenized Vaults'', https://eips.ethereum.org/EIPS/eip-4626.} It
also specifies rounding behavior:
\passthrough{\lstinline!convertToShares!} and
\passthrough{\lstinline!convertToAssets!} must round down.\footnote{Ethereum
  EIP-4626, \passthrough{\lstinline!convertToShares!} and
  \passthrough{\lstinline!convertToAssets!} rounding requirements,
  https://eips.ethereum.org/EIPS/eip-4626.} That rounding direction is
not just an implementation detail. It creates security properties and
integration risk.

Common share-price invariants:

\begin{lstlisting}
If no loss, fee, or donation occurs, share price should not decrease.

If a user deposits assets and immediately redeems the received shares
without intervening state changes, the loss must be bounded by documented
fees and rounding.

If totalSupply > 0, redeem(shares) should return assets proportional to
shares / totalSupply, subject to fees, liquidity, and rounding.

If totalSupply == 0, the initial exchange rate must not be manipulable in
a way that lets the first depositor confiscate later deposits.
\end{lstlisting}

The last property is where many real vault bugs live.

\subsection{Donation and Inflation Attacks Are Broken Share Invariants}

OpenZeppelin's ERC-4626 documentation describes the inflation attack: an
attacker deposits a small amount, donates assets directly to the vault,
shifts the exchange rate, and can make a later deposit mint zero or too
few shares due to rounding.\footnote{OpenZeppelin Contracts
  documentation, ``ERC4626'', security concern: inflation attack,
  https://docs.openzeppelin.com/contracts/5.x/erc4626.}

The broken invariant is not simply:

\begin{lstlisting}
sharePrice increased.
\end{lstlisting}

Share price may legitimately increase through yield or donations. The
broken property is closer to:

\begin{lstlisting}
A third party must not be able to manipulate the deposit exchange rate so that
a victim receives fewer shares than the protocol's slippage and rounding
policy allows, while the manipulator can recover the donated value through
their existing shares.
\end{lstlisting}

Attack shape:

\begin{lstlisting}
Initial state:
  totalAssets = 0
  totalShares = 0

Attacker:
  deposit 1 atomic unit, receive 1 share
  donate large amount directly to vault

Victim:
  deposit amount U
  receive 0 or too few shares because conversion rounds down

Attacker:
  redeem 1 share
  withdraw most or all assets
\end{lstlisting}

Mitigations change the invariant:

\begin{itemize}
\tightlist
\item
  Minimum initial liquidity.
\item
  Dead shares.
\item
  Virtual shares and virtual assets.
\item
  Deposit functions with minimum shares out.
\item
  Routers that protect first deposits.
\item
  Internal accounting that ignores direct donations.
\item
  Higher share precision.
\end{itemize}

Each mitigation has assumptions. Do not write ``prevent donation
attack.'' Write the exact post-mitigation property.

\subsection{AMM Invariants Are Pricing Rules, Not Magic Shields}

Uniswap documentation describes the constant product formula
\passthrough{\lstinline!x * y = k!}, where \passthrough{\lstinline!x!}
and \passthrough{\lstinline!y!} are reserves and
\passthrough{\lstinline!k!} is the invariant that stays constant or
increases after trades.\footnote{Uniswap Docs, ``How Uniswap Works'',
  https://developers.uniswap.org/docs/get-started/concepts/how-uniswap-works.}

That invariant is a pricing rule. It does not mean:

\begin{itemize}
\tightlist
\item
  The pool price is globally correct.
\item
  LPs cannot lose value.
\item
  Trades cannot be sandwiched.
\item
  Oracles derived from the pool are safe.
\item
  Fees cannot be misaccounted.
\item
  Reserves cannot be desynchronized from token balances.
\end{itemize}

AMM invariant examples:

\begin{lstlisting}
Swap:
  adjustedReserveXAfter * adjustedReserveYAfter >= reserveXBefore * reserveYBefore
  after accounting for fees.

Liquidity mint:
  minted LP shares must be proportional to contributed reserves,
  except for documented minimum liquidity and rounding.

Liquidity burn:
  burned LP shares must redeem proportional reserves,
  except for documented fees and rounding.

Reserve sync:
  stored reserves must match actual balances after operations that update reserves,
  and stale reserves must not be used as if they were fresh for security decisions.
\end{lstlisting}

The invariant tells you what kind of attack to search for:

\begin{itemize}
\tightlist
\item
  Can a token with transfer fees make stored reserves wrong?
\item
  Can a callback observe reserves before update?
\item
  Can a sync/skim mechanism move value unexpectedly?
\item
  Can a price oracle be manipulated within one block?
\item
  Can liquidity be minted at a favorable ratio because of rounding?
\end{itemize}

The formula is the beginning of review, not the end.

\subsection{Solvency Invariants Must Name Failure Modes}

Solvency can fail by:

\begin{itemize}
\tightlist
\item
  Theft.
\item
  Bad debt.
\item
  Mispriced collateral.
\item
  Strategy loss.
\item
  Rounding accumulation.
\item
  Untracked liabilities.
\item
  Frozen assets.
\item
  Withdrawal queue mismatch.
\item
  Bridge failure.
\item
  Governance seizure.
\item
  Liquidation malfunction.
\item
  Oracle outage.
\end{itemize}

Different failure modes require different invariants.

Example:

\begin{lstlisting}
Asset theft:
  Unauthorized accounts cannot reduce recoverableAssets.

Bad debt:
  After recognizing bad debt, supplier claims and reserve accounting must reflect
  the loss according to the documented socialization policy.

Frozen assets:
  Assets paused in a strategy cannot be counted as immediately withdrawable.

Oracle outage:
  Borrowing and liquidation behavior under stale oracle data must match the
  documented safety mode.
\end{lstlisting}

Solvency is not only ``assets \textgreater= liabilities.'' It is ``the
protocol's accounting of assets and liabilities matches reality under
each failure mode it claims to handle.''

\section{Temporal, Liveness, Exit, Withdrawal, and Recovery Invariants}

\subsection{Not Every Security Failure Is Theft}

A protocol can lose security without losing funds immediately.

Examples:

\begin{itemize}
\tightlist
\item
  Users cannot withdraw.
\item
  Messages never finalize.
\item
  A queue can be permanently blocked by one bad item.
\item
  A bridge pause has no exit path.
\item
  Liquidations cannot execute during market stress.
\item
  Oracle staleness freezes borrowing and repayment.
\item
  Governance cannot recover from a failed upgrade.
\item
  A rollup escape hatch exists on paper but cannot be used in practice.
\end{itemize}

These are liveness and recovery failures.

Safety property:

\begin{lstlisting}
No unauthorized withdrawal occurs.
\end{lstlisting}

Liveness property:

\begin{lstlisting}
An authorized withdrawal eventually becomes claimable if required conditions
remain satisfied and the system is not in a documented emergency state.
\end{lstlisting}

Both matter.

\subsection{Temporal Properties Need Clocks and Conditions}

Bad:

\begin{lstlisting}
Withdrawals should be processed quickly.
\end{lstlisting}

Better:

\begin{lstlisting}
If a withdrawal request is accepted at time t and the protocol remains unpaused,
then the request must become claimable no later than t + withdrawalDelay + maxEpochLength.
\end{lstlisting}

Better still:

\begin{lstlisting}
If a withdrawal request for amount a is accepted at epoch e,
and sufficient liquid assets or documented exit liquidity become available,
and no emergency pause is active,
then the request must become claimable by the protocol's stated processing rule
without later requests bypassing it except under documented priority rules.
\end{lstlisting}

Temporal properties always have conditions. Name them.

\subsection{Exit Invariants Are Security Boundaries}

An exit mechanism is a promise that users can leave.

Exit invariant:

\begin{lstlisting}
A user with a valid claim can eventually convert that claim into the underlying
asset or a documented settlement asset, unless a stated and bounded emergency
condition applies.
\end{lstlisting}

For bridges:

\begin{lstlisting}
A user who burns wrapped assets on the destination domain and submits a valid
message after source-domain finality must be able to unlock source assets
exactly once, unless the bridge is paused under documented recovery rules.
\end{lstlisting}

For rollups:

\begin{lstlisting}
A user must be able to force inclusion or exit through L1 under the rollup's
published censorship-resistance assumptions.
\end{lstlisting}

For vaults:

\begin{lstlisting}
A user who owns shares must be able to redeem according to liquidity, queue,
fee, and strategy-withdrawal rules, and the protocol must not silently convert
an immediate claim into an indefinite claim.
\end{lstlisting}

If an exit exists only when all operators are honest and online, then
operator honesty is part of the trust model. Write that down.

\subsection{Withdrawal Queues Need Ordering Invariants}

Queues are common:

\begin{itemize}
\tightlist
\item
  Staking withdrawals.
\item
  Vault redemptions.
\item
  Bridge exits.
\item
  Insurance claims.
\item
  Liquidation auctions.
\item
  Delayed settlement.
\end{itemize}

Queue invariant:

\begin{lstlisting}
If the design is FIFO, request i must not be processed after request j when
i was accepted before j, unless i is invalid, canceled, expired, or explicitly skipped
according to documented rules.
\end{lstlisting}

Capacity invariant:

\begin{lstlisting}
The sum of claimable withdrawals must not exceed liquid assets reserved for
the withdrawal batch.
\end{lstlisting}

Cancellation invariant:

\begin{lstlisting}
Canceling a withdrawal must restore exactly the user's unclaimed rights
minus documented fees, and must not let the user double claim.
\end{lstlisting}

Griefing invariant:

\begin{lstlisting}
No user can submit a request that permanently blocks unrelated valid requests
without bearing a cost proportional to the damage or without an administrative
recovery path.
\end{lstlisting}

Queues convert accounting bugs into time bugs. Review both.

\subsection{Pause and Recovery Controls Need Invariants Too}

Emergency controls are security-critical because they intentionally
break normal flow.

Pause invariant:

\begin{lstlisting}
When paused:
  prohibited actions must revert,
  permitted recovery actions must remain available,
  user claim accounting must not be corrupted,
  and the pause authority must not gain extra withdrawal or upgrade power
  unless explicitly documented.
\end{lstlisting}

Unpause invariant:

\begin{lstlisting}
Unpause must require the documented authority and must not skip required
post-incident checks, timelocks, or governance confirmations.
\end{lstlisting}

Recovery invariant:

\begin{lstlisting}
A recovery action must move the system from a known bad or degraded state
to a documented safe state without increasing total user claims, reducing
user claims arbitrarily, or granting permanent emergency authority.
\end{lstlisting}

Bad recovery controls create second exploits after the first incident.

\subsection{Liveness Invariants Are Hard to Test}

Unit tests are often bad at liveness because liveness is about time,
eventuality, and environment behavior.

Testing strategies:

\begin{itemize}
\tightlist
\item
  Model time explicitly.
\item
  Advance blocks, slots, or epochs.
\item
  Simulate paused and unpaused states.
\item
  Simulate empty and full queues.
\item
  Simulate failed external calls.
\item
  Simulate partial liquidity.
\item
  Simulate oracle downtime.
\item
  Simulate operator absence.
\item
  Test that recovery paths remain callable.
\end{itemize}

Formal verification can express some temporal properties, but many
blockchain teams rely on stateful tests, model checking, simulations,
and manual reasoning. The key is not tool purity. The key is that the
property is written.

\section{Exploitability: Capital, Timing, Liquidity, Reliability, and Profit}

\subsection{Vulnerability Is Not the Same as Exploit}

A vulnerability is a weakness. An exploit is a feasible strategy that
uses the weakness to achieve impact.

Example vulnerability:

\begin{lstlisting}
The vault's first-deposit exchange rate can be manipulated by donation.
\end{lstlisting}

Exploitability questions:

\begin{itemize}
\tightlist
\item
  Is the vault empty or can it become empty?
\item
  Can the attacker front-run the victim's deposit?
\item
  Can the attacker donate directly to the vault?
\item
  Does rounding mint zero or too few shares?
\item
  Does the victim transaction set a minimum share output?
\item
  Can the attacker redeem after the victim deposit?
\item
  How much capital is required?
\item
  Is the attack profitable after gas and opportunity cost?
\item
  Can the attack be repeated?
\item
  How many users or assets are exposed?
\end{itemize}

The same code flaw can be catastrophic in one deployment and low
severity in another because the exploitability conditions differ.

\subsection{Exploitability Has Preconditions}

Always write preconditions.

Template:

\begin{lstlisting}
Threat:
  Attacker violates invariant P.

Preconditions:
  Initial state:
  Required attacker capital:
  Required permissions:
  Required timing or ordering:
  Required liquidity:
  Required external dependency behavior:
  Required victim behavior:
  Required protocol configuration:

Attack path:
  Step 1:
  Step 2:
  Step 3:

Impact:
  Direct profit:
  Funds at risk:
  Denial or griefing:
  Governance or control impact:
  Recovery difficulty:

Reliability:
  Deterministic or probabilistic:
  Competition:
  Monitoring/detection:
  Repeatability:
\end{lstlisting}

If you cannot fill the preconditions, you do not yet understand
exploitability.

\subsection{Capital Requirement}

Capital matters, but not simplistically.

Questions:

\begin{itemize}
\tightlist
\item
  How much capital must be owned?
\item
  Can it be flash-borrowed?
\item
  Must it be held across blocks?
\item
  Is it at market risk?
\item
  Is it locked during a challenge window?
\item
  Is it consumed as a cost or recoverable?
\item
  Does the attack require scarce collateral?
\item
  Does slippage scale with attack size?
\item
  Does the attacker need governance voting power?
\end{itemize}

Examples:

\begin{lstlisting}
Flash-loan feasible:
  Attacker needs large temporary liquidity within one transaction.

Capital locked:
  Attacker must stake assets for 7 days and faces slashing risk.

Unborrowable capital:
  Attacker needs governance voting power at a historical snapshot.

Inventory risk:
  Attacker must hold volatile collateral while manipulating a market over time.
\end{lstlisting}

``Requires a lot of capital'' is not a defense unless the amount,
source, duration, and risk are specified.

\subsection{Timing and Ordering}

Blockchain exploits often depend on ordering.

Questions:

\begin{itemize}
\tightlist
\item
  Must the attacker act before a victim transaction?
\item
  Must the attacker act after an oracle update?
\item
  Must several actions occur in one transaction?
\item
  Can a private mempool or builder relationship help?
\item
  Can a sequencer censor or reorder?
\item
  Can a keeper race the attacker?
\item
  Does the attack survive reorgs or finality delays?
\item
  Is there a challenge period?
\item
  Does the exploit require a stale price window?
\end{itemize}

Ordering-dependent exploit:

\begin{lstlisting}
Attacker front-runs first deposit, manipulates share price, victim deposits,
attacker back-runs redemption.
\end{lstlisting}

Temporal exploit:

\begin{lstlisting}
Attacker borrows against collateral immediately after a manipulated oracle update
and before the price reverts.
\end{lstlisting}

Liveness exploit:

\begin{lstlisting}
Attacker fills a withdrawal queue with dust requests so large withdrawals cannot
be processed before a deadline.
\end{lstlisting}

Write ordering as part of the exploit, not as a footnote.

\subsection{Liquidity and Market Depth}

Economic attacks need markets.

Questions:

\begin{itemize}
\tightlist
\item
  Is there enough liquidity to manipulate the price?
\item
  Is there enough liquidity to exit the profit?
\item
  Does the manipulation cost exceed the extracted value?
\item
  Is the oracle based on spot price, TWAP, median, or signed feeds?
\item
  How long must manipulation persist?
\item
  Can arbitrage restore the price before the exploit completes?
\item
  Does the attacker need to trade through fee tiers or concentrated
  ranges?
\item
  Are there liquidity cliffs?
\end{itemize}

Invariant:

\begin{lstlisting}
The oracle price used for collateral checks must not be manipulable within the
capital and time bounds available to a rational attacker for profit greater
than manipulation cost.
\end{lstlisting}

That is not directly a simple code assertion. It may require simulation,
market analysis, and risk parameter review. But it is still a security
property.

\subsection{Reliability and Repeatability}

A one-shot exploit with uncertain success is different from a
deterministic drain.

Reliability questions:

\begin{itemize}
\tightlist
\item
  Does the exploit always work if preconditions hold?
\item
  Does it depend on miner, validator, sequencer, or builder cooperation?
\item
  Does it depend on victim behavior?
\item
  Can it be simulated exactly on a fork?
\item
  Can failed attempts revert without cost?
\item
  Can competitors copy the transaction?
\item
  Does the protocol detect or pause after the first attempt?
\item
  Can the attacker repeat across markets, assets, chains, or users?
\end{itemize}

Repeatability often determines severity.

\begin{lstlisting}
Single victim:
  Profit limited to one pending transaction.

Per-market:
  Attack repeats across all markets with same accounting flaw.

Global drain:
  Attack can drain all assets controlled by the protocol.

Permanent control:
  Attack grants upgrade, mint, or governance authority.
\end{lstlisting}

\subsection{Profit Is Not the Only Impact}

Some attacks are not profitable thefts.

They may cause:

\begin{itemize}
\tightlist
\item
  Permanent fund lock.
\item
  Forced liquidation.
\item
  Bad debt.
\item
  Governance takeover.
\item
  Oracle corruption.
\item
  Chain halt.
\item
  Bridge halt.
\item
  Censorship.
\item
  User deanonymization.
\item
  Loss of accounting integrity.
\item
  Loss of upgrade control.
\item
  Griefing at low cost.
\end{itemize}

Security severity should not require attacker profit. A cheap griefing
attack that freezes a billion-dollar withdrawal queue is serious even if
the attacker cannot directly steal.

\subsection{``Not Exploitable'' Must Be Proven Narrowly}

Bad statement:

\begin{lstlisting}
This is not exploitable.
\end{lstlisting}

Better:

\begin{lstlisting}
Under the current deployment, the attack requires the vault to be empty.
The vault has non-burnable dead shares and a minimum initial liquidity position.
A direct donation can change share price, but deposits route through a function
that enforces minimum shares out. Therefore the specific zero-share first-deposit
attack is not feasible unless those conditions change.
\end{lstlisting}

This statement is useful because it names:

\begin{itemize}
\tightlist
\item
  Deployment condition.
\item
  Mitigations.
\item
  Remaining effect.
\item
  Attack variant ruled out.
\item
  Configuration sensitivity.
\end{itemize}

``Not exploitable'' without scope is usually false confidence.

\subsection{Exploitability Worksheet}

Use this worksheet for every serious invariant violation.

\begin{lstlisting}
Invariant:

Violation:

Smallest counterexample:

Required initial state:

Required attacker capabilities:

Required victim or keeper behavior:

Required external dependency behavior:

Capital needed:

Time/order needed:

Liquidity/depth needed:

Gas/compute/storage constraints:

Profit or impact:

Repeatability:

Detection likelihood:

Mitigations:

Residual risk:
\end{lstlisting}

If the counterexample exists but exploitability is unclear, say so. That
is not a failure. It is honest security analysis.

\section{Worked Example: From Protocol Intent to Attack Hypotheses}

\subsection{Protocol Description}

Consider a simplified tokenized vault:

\begin{lstlisting}
Users deposit an ERC-20 asset and receive vault shares.
Users redeem shares for a proportional amount of the vault's assets.
The vault uses asset.balanceOf(vault) as total assets.
When totalShares == 0, initial deposits mint shares 1:1 with assets.
When totalShares > 0:
  sharesOut = assetsIn * totalShares / asset.balanceOf(vault)
rounded down.
\end{lstlisting}

The intended security story:

\begin{itemize}
\tightlist
\item
  Shares represent proportional claims on vault assets.
\item
  Depositors should not be diluted unexpectedly.
\item
  Rounding should be small.
\item
  A user who deposits and later redeems should receive their
  proportional claim.
\item
  Anyone may deposit.
\item
  Anyone may transfer assets directly to the vault, because ERC-20
  transfers cannot be prevented in general.
\end{itemize}

Now write invariants.

\subsection{Step 1: Asset and Share Conservation}

Basic accounting:

\begin{lstlisting}
totalShares == sum(userShares)
\end{lstlisting}

Backing:

\begin{lstlisting}
asset.balanceOf(vault) >= assetsRedeemableByAllShares
\end{lstlisting}

But this is not enough. If
\passthrough{\lstinline!assetsRedeemableByAllShares!} is computed from
the same manipulated \passthrough{\lstinline!asset.balanceOf(vault)!},
this invariant may be true during the attack.

The deeper claim is user fairness under deposit pricing.

\subsection{Step 2: Deposit Fairness Invariant}

For normal deposits:

\begin{lstlisting}
If user deposits amount a through the official deposit function,
and no fee or loss condition applies,
then shares minted must represent a proportional claim on the deposited assets,
with rounding loss bounded by the protocol's documented tolerance.
\end{lstlisting}

This is still too vague. Add attack context:

\begin{lstlisting}
For any attacker-controlled direct donation before a victim deposit,
the attacker must not be able to recover the donation plus a material portion
of the victim deposit solely because the donation changed the share conversion rate.
\end{lstlisting}

Now we have a target.

\subsection{Step 3: Empty-Vault Attack Hypothesis}

Hypothesis:

\begin{lstlisting}
If the vault is empty, an attacker can become the only shareholder with a tiny
deposit, donate assets directly to inflate asset.balanceOf(vault), and cause a
victim deposit to mint zero or too few shares because sharesOut rounds down.
\end{lstlisting}

Preconditions:

\begin{lstlisting}
Vault state:
  totalShares == 0 or very small

Attacker capabilities:
  Can deposit before victim
  Can transfer asset directly to vault
  Can redeem shares after victim deposit

Victim behavior:
  Victim deposit does not enforce minimum shares out
  Victim transaction can be observed or anticipated

Token behavior:
  Direct transfers increase asset.balanceOf(vault)

Math:
  sharesOut rounds down

Capital:
  Attacker can temporarily supply donation amount comparable to victim deposit
\end{lstlisting}

Attack path:

\begin{lstlisting}
1. Attacker deposits 1 atomic unit.
2. Vault mints 1 share to attacker.
3. Attacker transfers D assets directly to the vault.
4. Victim deposits U assets.
5. sharesOut = U * 1 / (1 + D + maybe U depending on implementation)
6. If sharesOut rounds to 0 or too few shares, victim receives little or no claim.
7. Attacker redeems 1 share for most vault assets.
\end{lstlisting}

Invariant violation:

\begin{lstlisting}
Victim's deposit minted shares with rounding loss far above the documented
or economically acceptable bound, and the loss was captured by the attacker.
\end{lstlisting}

Impact:

\begin{lstlisting}
The attacker can steal a victim's first deposit or a material fraction of it.
\end{lstlisting}

This is not hypothetical as a class of bug. It is the ERC-4626 inflation
attack pattern discussed in the OpenZeppelin documentation.\footnote{OpenZeppelin
  Contracts documentation, ``ERC4626'', security concern: inflation
  attack, https://docs.openzeppelin.com/contracts/5.x/erc4626.}

\subsection{Step 4: Exploitability Evaluation}

Capital:

\begin{lstlisting}
The attacker needs enough assets to donate so the victim's share output rounds
to the target level. The donation may be mostly recoverable if the attacker is
the dominant shareholder.
\end{lstlisting}

Timing:

\begin{lstlisting}
The attacker benefits from front-running or otherwise preceding the victim's
deposit while the vault has low share supply.
\end{lstlisting}

Liquidity:

\begin{lstlisting}
The attack uses the vault asset itself, not necessarily external market liquidity.
If capital can be flash-borrowed and the deposit/redeem sequence fits in one
transaction or bundle, capital constraints may be lower.
\end{lstlisting}

Reliability:

\begin{lstlisting}
If the vault is empty and deposits lack minimum-share protection, the attack is
highly reliable. If the vault already has large share supply, the attack may be
uneconomic.
\end{lstlisting}

Repeatability:

\begin{lstlisting}
Usually strongest at initialization or after a vault is emptied. It may repeat
across newly deployed vaults, markets, or assets.
\end{lstlisting}

Severity:

\begin{lstlisting}
High if users can make unprotected first deposits and assets are valuable.
Lower if deployment process seeds the vault safely, users deposit through a
router with slippage protection, or virtual offsets make the attack uneconomic.
\end{lstlisting}

\subsection{Step 5: Mitigation Invariants}

Virtual shares/assets:

\begin{lstlisting}
The initial exchange rate must be anchored so that a direct donation cannot
make victim rounding loss exceed the documented bound at economically rational
capital levels.
\end{lstlisting}

Minimum shares out:

\begin{lstlisting}
A deposit must revert if sharesOut < minSharesOut supplied by the depositor
or caller.
\end{lstlisting}

Seed liquidity:

\begin{lstlisting}
Before public deposits are enabled, the vault must have initial liquidity and
share supply sufficient to bound rounding loss for expected deposit sizes.
\end{lstlisting}

Internal accounting:

\begin{lstlisting}
Direct token transfers must not affect deposit pricing unless the protocol
explicitly treats donations as yield and has bounded donation-capture risk.
\end{lstlisting}

Monitoring:

\begin{lstlisting}
Alert if a vault with low totalShares receives a direct asset transfer before
or near a public deposit.
\end{lstlisting}

Residual risk:

\begin{lstlisting}
Even with mitigations, integrators using preview functions as hard price oracles
may be exposed if they ignore slippage, rounding, or manipulable on-chain state.
\end{lstlisting}

This is the full workflow:

\begin{lstlisting}
Intent -> invariant -> counterexample -> exploitability -> mitigation invariant -> tests and monitoring
\end{lstlisting}

\section{Practical Discipline: Building an Invariant Catalogue}

For every protocol, maintain an invariant catalogue. It should be short
enough to use and precise enough to guide review.

Template:

\begin{lstlisting}
Invariant ID:

Name:

Type:
  state / transition / temporal / economic / authorization / cross-domain

Statement:

Scope:
  contracts, assets, users, domains, roles

Assumptions:
  oracle, token behavior, finality, governance, operators, liquidity

Allowed exceptions:

Impact if violated:

Likely attack paths:

Tests:
  unit:
  fuzz/invariant:
  scenario:
  fork/simulation:
  formal/spec:

Monitoring:

Residual risk:
\end{lstlisting}

Example:

\begin{lstlisting}
Invariant ID:
  VAULT-ROUNDING-001

Name:
  Deposit rounding loss is bounded

Type:
  transition / economic

Statement:
  For any successful deposit of assets a, shares minted to receiver must be
  at least the minimum output specified by the caller and must not create
  attacker-capturable rounding loss above the configured tolerance.

Scope:
  All vault deposit and mint paths, including router and direct calls.

Assumptions:
  Asset token does not reenter during transfer.
  Asset decimals are configured correctly.
  Direct donations are either ignored by internal accounting or made uneconomic
  by virtual offset.

Allowed exceptions:
  Documented deposit fee.
  Revert if minSharesOut is not met.

Impact if violated:
  User deposit can be partially or fully captured by existing shareholders.

Likely attack paths:
  First-deposit donation.
  Low-supply rounding manipulation.
  Fee-on-transfer mismatch.

Tests:
  Unit test empty vault first deposit.
  Stateful fuzz deposits, donations, redeems, low share supply.
  Fork test against supported assets with real decimals.

Monitoring:
  Alert on direct transfers into low-supply vaults.
\end{lstlisting}

This kind of catalogue becomes the bridge to later chapters on testing,
formal verification, and review workflow.

\section{Review Questions}

\begin{enumerate}
\def\labelenumi{\arabic{enumi}.}
\tightlist
\item
  Why is ``users should not lose funds'' not a useful invariant?
\item
  What is the difference between an invariant over all states and an invariant over reachable states?
\item
  Give one example of a state invariant, a transition invariant, a temporal property, and an economic assumption.
\item
  Why can a true implementation invariant still fail to capture protocol intent?
\item
  What assumptions would you attach to a lending collateralization invariant?
\item
  Why is \passthrough{\lstinline!totalSupply == sum(balances)!} insufficient for a wrapped asset?
\item
  What must be specified before claiming that a bug is ``not exploitable''?
\item
  How do capital, timing, liquidity, ordering control, and reliability change exploit severity?
\end{enumerate}

\section{Applied Exercises}

\subsection{Exercise 1: Write Invariants from Intent}

Protocol intent:

\begin{lstlisting}
Users deposit Token A into a vault and receive shares.
The vault deploys Token A into a strategy.
Users can request withdrawal.
Withdrawals are processed once per epoch.
The strategist can move funds between approved strategies.
Governance can add strategies after a 2-day timelock.
\end{lstlisting}

Write:

\begin{itemize}
\tightlist
\item
  three conservation invariants;
\item
  two authorization invariants;
\item
  two temporal or withdrawal properties;
\item
  one recovery invariant for strategy failure.
\end{itemize}

Expected output: an invariant catalogue with assumptions for each invariant.

\subsection{Exercise 2: Find the Missing Assumption}

Invariant:

\begin{lstlisting}
Borrowing is allowed only if debtValue <= collateralValue * collateralFactor.
\end{lstlisting}

List at least six assumptions hidden inside this invariant. Consider oracle source, freshness, decimal scaling, liquidity, interest accrual, token transfer behavior, governance parameter changes, sequencer availability, and liquidation liveness.

Expected output: an assumption table with a failure mode for each assumption.

\subsection{Exercise 3: Exploitability Worksheet}

Threat:

\begin{lstlisting}
An attacker can manipulate an AMM spot price used by a lending protocol oracle.
\end{lstlisting}

Fill out:

\begin{itemize}
\tightlist
\item
  broken invariant;
\item
  preconditions;
\item
  required capital;
\item
  required timing;
\item
  liquidity assumptions;
\item
  attack path;
\item
  profit calculation inputs;
\item
  reliability;
\item
  mitigations;
\item
  residual risk.
\end{itemize}

Expected output: an exploitability worksheet and a conclusion: feasible, infeasible, or uncertain.

\subsection{Exercise 4: Counterexample to Finding}

Counterexample:

\begin{lstlisting}
1. Vault is empty.
2. Attacker deposits 1 unit and receives 1 share.
3. Attacker donates 100,000 units directly to the vault.
4. Victim deposits 50,000 units and receives 0 shares.
5. Attacker redeems 1 share for almost all assets.
\end{lstlisting}

Write:

\begin{itemize}
\tightlist
\item
  root cause;
\item
  broken invariant;
\item
  preconditions;
\item
  impact;
\item
  exploitability assessment;
\item
  recommended mitigation;
\item
  tests that should be added;
\item
  residual risk after mitigation.
\end{itemize}

Expected output: a concise finding draft.

\section{Key Takeaways}

\begin{itemize}
\tightlist
\item
  An invariant is a property that should hold over reachable states,
  transitions, or histories.
\item
  Many serious blockchain exploits are valid transaction sequences that
  violate intended invariants.
\item
  Useful invariants are specific, falsifiable, assumption-aware,
  measurable, and linked to impact.
\item
  Conservation invariants ask whether claims exceed backing.
\item
  Authorization invariants ask who can cause protected state changes,
  through which paths, under what conditions.
\item
  Solvency invariants must include recoverability, liabilities,
  reserves, prices, and failure modes.
\item
  Share-price invariants are vulnerable to rounding, direct donations,
  first-deposit conditions, and integration assumptions.
\item
  Temporal and liveness properties cover withdrawals, queues, exits,
  pauses, and recovery.
\item
  Exploitability analysis turns a counterexample into a realistic attack
  assessment.
\item
  ``Not exploitable'' is meaningful only when scoped to concrete
  preconditions, mitigations, deployment facts, and attacker
  capabilities.
\item
  Tools can test or prove properties, but humans must choose the right
  properties.
\end{itemize}

\section{Part I Handoff}

Part I ends with an Integrated Lab rather than a separate capstone. The lab asks the reader to combine the first seven chapters into one security-model review package.

The required artifact is:

\begin{lstlisting}
Security Model Review Package:
  state-machine diagram
  transition table
  asset inventory
  authority map
  trust-boundary map
  transaction analysis
  authorization-flow map
  threat model
  invariant catalogue
  exploitability memo
\end{lstlisting}

Use \passthrough{\lstinline!part-01-integrated-lab.tex!} as the working file for that review.

Part II then applies the same method below the application layer. Consensus, fork choice, finality, peer-to-peer networks, block production, MEV, incentives, clients, forks, upgrades, and chain incidents determine which histories become real, which transactions can be included or censored, which assumptions applications inherit, and which failures no smart contract can fix locally.

\section{Further Reading}

\begin{itemize}
\tightlist
\item
  Leslie Lamport,
  \href{https://lamport.azurewebsites.net/pubs/state-machine.pdf}{``Computation
  and State Machines''}. Foundation for state-machine and invariant
  reasoning.
\item
  Bowen Alpern and Fred B. Schneider,
  \href{https://www.cs.cornell.edu/fbs/publications/DefLiveness.pdf}{``Defining
  Liveness''}. Classical background on safety and liveness properties.
\item
  Foundry,
  \href{https://www.getfoundry.sh/guides/invariant-testing}{``Invariant
  Testing''}. Practical invariant testing guidance for Solidity
  projects.
\item
  Trail of Bits,
  \href{https://secure-contracts.com/program-analysis/echidna/introduction/how-to-test-a-property.html}{``Testing
  a Property with Echidna''} and
  \href{https://secure-contracts.com/program-analysis/echidna/basic/testing-modes.html}{``How
  to Select the Most Suitable Testing Mode''}. Property-based and
  stateful smart-contract fuzzing guidance.
\item
  Solidity documentation,
  \href{https://docs.soliditylang.org/en/latest/smtchecker.html}{``SMTChecker
  and Formal Verification''}. Official Solidity guidance on assertions,
  state properties, counterexamples, and specification limits.
\item
  Certora,
  \href{https://docs.certora.com/en/latest/docs/cvl/invariants.html}{``Invariants''}.
  Practical explanation of invariant preservation, induction, and
  pitfalls in formal specifications.
\item
  Ethereum EIP-4626,
  \href{https://eips.ethereum.org/EIPS/eip-4626}{``Tokenized Vaults''}.
  Standardized vault share/accounting interface and rounding
  requirements.
\item
  OpenZeppelin,
  \href{https://docs.openzeppelin.com/contracts/5.x/erc4626}{``ERC4626''}.
  Includes an explanation of the ERC-4626 inflation attack and
  virtual-offset defense.
\item
  Uniswap,
  \href{https://developers.uniswap.org/docs/get-started/concepts/how-uniswap-works}{``How
  Uniswap Works''}. Useful reference for AMM reserve and
  constant-product invariant intuition.
\item
  Aave, \href{https://aave.com/help/borrowing/liquidations}{``Health
  Factor \& Liquidations''}. Practical description of lending health
  factor and liquidation conditions.
\item
  Compound v2,
  \href{https://docs.compound.finance/v2/comptroller/}{``Comptroller''}.
  Versioned reference for risk management, collateral factors,
  liquidity, and shortfall.
\end{itemize}
